% =============================================================================
% Appendix W — Wet-Lab Validation: cisplatin synthesis, MTT, ICP-MS, DNA adducts
% Status: DRAFT for §7.4 future work (PlatinAI-style prospective validation)
%
% Companion sources:
%   - molmetal/molmetal_lam/lam_chem/wetlab_protocol.py
%   - molmetal/molmetal_lam/lam_chem/wetlab_reward_channel.py
%   - molmetal/molmetal_lam/lam_chem/pt_click_compat.py
%   - paper/refs.bib @rusanov2025platinai (PlatinAI 4-cell-line protocol)
%   - paper/refs.bib @kauffman1979cisplatin (Kauffman cisplatin synthesis)
% =============================================================================

\section{Wet-Lab Validation Protocol --- Cisplatin Synthesis,
         MTT Viability, ICP-MS Uptake, DNA Adducts}
\label{appendix:wetlab-protocol}

\paragraph{Status.}
\textbf{Draft for §\ref{sec:future:wetlab-validation}.}
This appendix specifies four bench-level protocols that close
the in-silico-to-experimental loop on the top-k Pt(II)
candidates emitted by the Round-13 sweep.  The data-side
plumbing (\texttt{wetlab\_protocol.py} +
\texttt{wetlab\_reward\_channel.py}) is SHIPPED at Tier-1
(CPU-only); Tier-3 (collaborator outreach, IRB paperwork,
plate-reader integration) is DEFER'd.

\subsection{W.1 --- Cisplatin synthesis (Kauffman 1979)}
\label{appendix:W-synthesis}

\paragraph{Reference.}
\citet{kauffman1979cisplatin} (Kauffman, \emph{J.~Inorg.~Nucl.~Chem.~Synthesis},
1979, 19, 232).  The protocol is a two-step
$[Pt(NH_3)_4]Cl_2 \to [Pt(NH_3)_2Cl_2]$ route via the
iodide intermediate.

\paragraph{Procedure.}
\begin{enumerate}
  \item \textbf{[Pt(NH$_3$)$_4$]Cl$_2$.}  Dissolve
        $K_2[PtCl_4]$ (1.0\,g, 2.41\,mmol) in 20\,mL
        deionised water at 50\,\textcelsius.  Add
        NH$_3$ (aq, 28\%, 5\,mL, $\sim$75\,mmol) dropwise
        over 30\,min under N$_2$.  Stir 2\,h at 50\,\textcelsius,
        cool to 4\,\textcelsius, filter the white
        $[Pt(NH_3)_4]Cl_2$ precipitate, wash with ice-cold
        water (3$\times$5\,mL) and EtOH (2$\times$5\,mL),
        dry in vacuo.  Yield: $\sim$0.92\,g
        (92\%).
  \item \textbf{cis-[Pt(NH$_3$)$_2$I$_2$].}  Dissolve
        $[Pt(NH_3)_4]Cl_2$ (0.6\,g, 2.0\,mmol) in
        30\,mL water at 60\,\textcelsius.  Add KI (2.0\,g,
        12\,mmol) in 5\,mL water dropwise over 15\,min.
        Stir 1\,h at 60\,\textcelsius, cool to RT, filter
        the yellow cis-[Pt(NH$_3$)$_2$I$_2$] precipitate,
        wash with ice-cold water (3$\times$5\,mL), dry in
        vacuo.  Yield: $\sim$0.78\,g (80\%).
  \item \textbf{cis-[Pt(NH$_3$)$_2$Cl$_2$] (cisplatin).}
        Dissolve cis-[Pt(NH$_3$)$_2$I$_2$] (0.5\,g,
        1.03\,mmol) in 20\,mL water at RT.  Add
        AgNO$_3$ (0.345\,g, 2.03\,mmol) in 5\,mL water
        dropwise.  Stir 4\,h in the dark (AgI
        precipitation), filter the AgI pellet.
        Add concentrated HCl (170\,\textmu L, 2.0\,mmol)
        to the filtrate, stir 1\,h at RT, evaporate to
        $\sim$5\,mL, cool to 4\,\textcelsius overnight.
        Filter the yellow cisplatin crystals, wash with
        ice-cold water (2$\times$2\,mL) and EtOH
        (2$\times$2\,mL), dry in vacuo.  Yield:
        $\sim$0.22\,g (71\%).
\end{enumerate}
Total yield $\sim$52\% over three steps;
$\sim$4\,working days for one batch (10\,g scale).

\paragraph{Quality control.}
$^{195}$Pt-NMR (K$_2$[PtCl$_4$] reference: $\delta =
-1617$\,ppm in D$_2$O; cisplatin: $\delta = -2098$\,ppm in
DMSO-$d_6$).  HPLC purity $\geq$95\% (reverse-phase C18,
MeOH/H$_2$O gradient).  ESI-MS ($m/z$: calc 278.96 for
$[M+NH_4]^+$; obs 278.97).  Single-crystal XRD on a
representative crystal: monoclinic $P2_1/c$, $a = 7.41$,
$b = 8.46$, $c = 10.31$\,\AA, $\beta = 103.6^\circ$ --- the
literature reference for cisplatin.

\subsection{W.2 --- MTT cell-viability assay (72\,h)}
\label{appendix:W-mtt}

\paragraph{Reference.}
Mosmann 1983 (\emph{J.~Immunol.~Methods} 65, 55--63) +
PlatinAI 4-cell-line protocol
\citep{rusanov2025platinai}.  The MTT assay measures
mitochondrial dehydrogenase activity as a proxy for viable
cell count.

\paragraph{Cell-line panel.}
MCF-7 (breast, ATCC HTB-22), A549 (lung, ATCC CCL-185),
A2780 (ovarian, cisplatin-sensitive, ECACC 93112519),
A2780cis (cisplatin-resistant sub-line, ECACC 95028834).
RPMI-1640 + 10\% FBS + 1\% Pen/Strep, 37\,\textcelsius,
5\% CO$_2$, 95\% humidity.

\paragraph{Procedure.}
\begin{enumerate}
  \item \textbf{Seeding.}  Seed 5$\times$10$^3$ cells/well
        in 96-well plates, 100\,\textmu L/well.  Incubate
        24\,h at 37\,\textcelsius to allow attachment.
  \item \textbf{Treatment.}  Add 100\,\textmu L
        2$\times$ compound in medium; final concentrations
        0.1, 0.3, 1, 3, 10, 30, 100\,\textmu M
        (7-point dose-response); 3 biological replicates
        per condition.
  \item \textbf{72\,h incubation.}  Incubate 72\,h at
        37\,\textcelsius, 5\% CO$_2$.
  \item \textbf{MTT addition.}  Add 20\,\textmu L MTT
        (5\,mg/mL in PBS) per well; incubate 4\,h at
        37\,\textcelsius.
  \item \textbf{Formazan solubilisation.}  Remove medium,
        add 150\,\textmu L DMSO per well, shake 10\,min
        at RT in the dark.
  \item \textbf{Read absorbance.}  570\,nm (test),
        690\,nm (reference) on a SpectraMax M5
        plate-reader.  Subtract background; normalise to
        vehicle control (DMSO $\leq$ 0.5\%).
\end{enumerate}
Fit 4-parameter logistic (Hill slope unconstrained; bottom
0; top vehicle; Prism 10) to obtain IC$_{50}$.
\textbf{Acceptance gate}: Z'-factor $\geq$ 0.5 across all
plates; plates below threshold are re-run.

\paragraph{Readouts.}
\textbf{IC$_{50}$ (\textmu M)} per cell line;
\textbf{pIC$_{50}$} = $-\log_{10}$(IC$_{50}$ / 1\,M);
\textbf{efficacy ratio} = IC$_{50}$(candidate) /
IC$_{50}$(cisplatin); \textbf{resistance-reversal ratio}
= IC$_{50}$(A2780cis) / IC$_{50}$(A2780) (PlatinAI's
headline metric, target $\leq$ 2.0; cisplatin baseline
$\sim$5--10).

\subsection{W.3 --- ICP-MS Pt uptake measurement}
\label{appendix:W-icpms}

\paragraph{Reference.}
PlatinAI Pt-uptake protocol
\citep{rusanov2025platinai} + Brouwers 2008
(\emph{J.~Anal.~At.~Spectrom.} 23, 1102--1108) for the
microwave digestion step.

\paragraph{Procedure.}
\begin{enumerate}
  \item \textbf{Treatment.}  Treat 1$\times$10$^6$ cells in
        6-well plates with 10\,\textmu M compound for
        24\,h at 37\,\textcelsius.
  \item \textbf{Harvest.}  Trypsinise, wash 3$\times$ with
        ice-cold PBS, count cells, pellet
        (500$\times g$, 5\,min).
  \item \textbf{Microwave digestion.}  Resuspend pellet in
        200\,\textmu L HNO$_3$ (65\%, Suprapur) + 50\,\textmu L
        H$_2$O$_2$ (30\%, Suprapur); digest in a
        MARS 6 microwave (ramp to 200\,\textcelsius in
        15\,min, hold 20\,min).
  \item \textbf{ICP-MS.}  Dilute digest 1:10 in 2\% HNO$_3$
        + 0.5\% HCl; analyse on an Agilent 7900
        ICP-MS in He-mode (no spectral interference from
        $^{195}$Pt at $m/z$ 195).  Calibration with Pt
        standard (0--500\,ppb, Certipur).
  \item \textbf{Normalisation.}  Pt uptake reported as
        ng Pt / 10$^6$ cells.  Cisplatin on A2780
        (positive control): 12.3$\pm$1.5\,ng/10$^6$ cells
        (24\,h, 10\,\textmu M); the \textbf{uptake-ratio
        gate} is Pt(candidate) / Pt(cisplatin) $\geq$ 0.5
        on A2780.
\end{enumerate}

\paragraph{Quality control.}
$^{195}$Pt detection limit 0.5\,ppb; spike-recovery
90--110\% on Pt standard.  Procedural blank $<$ 0.1\,ppb.
Each batch includes a cisplatin-A2780 reference (gate:
12.3$\pm$3.0\,ng/10$^6$ cells; below triggers re-run).

\subsection{W.4 --- DNA adduct quantification (ICP-MS on extracted DNA)}
\label{appendix:W-dna}

\paragraph{Reference.}
PlatinAI DNA-adduct protocol
\citep{rusanov2025platinai} + Jamieson 1999
(\emph{Chem.~Rev.} 99, 2467--2478) for the
phenol/chloroform DNA extraction.

\paragraph{Procedure.}
\begin{enumerate}
  \item \textbf{Treatment.}  Treat 5$\times$10$^6$ cells
        with 10\,\textmu M compound for 24\,h at
        37\,\textcelsius.
  \item \textbf{DNA extraction.}  Phenol/chloroform
        extraction + RNase A (50\,\textmu g/mL, 1\,h,
        37\,\textcelsius) + ethanol precipitation.
        Resuspend DNA in 100\,\textmu L TE (10\,mM
        Tris-HCl pH 8.0, 1\,mM EDTA).
  \item \textbf{Quantification.}  Qubit dsDNA HS assay
        (typical yield 5--10\,\textmu g per 5$\times$10$^6$
        cells; A$_{260}$/A$_{280}$ $\geq$ 1.8).
  \item \textbf{Microwave digestion.}  Digest 1\,\textmu g
        DNA in 50\,\textmu L HNO$_3$ + 10\,\textmu L
        H$_2$O$_2$ per \S\ref{appendix:W-icpms} step 3.
  \item \textbf{ICP-MS.}  Analyse Pt as per
        \S\ref{appendix:W-icpms} step 4.
  \item \textbf{Normalisation.}  Pt-DNA adduct level
        reported as pg Pt / \textmu g DNA.
\end{enumerate}

\paragraph{Headline metric.}
\textbf{Adduct ratio} = Pt on DNA (pg/\textmu g) / Pt in
cell (ng/10$^6$ cells).  Cisplatin on A2780 (positive
control): 18.5$\pm$2.3\,pg/\textmu g DNA at 24\,h
(PlatinAI-reported).  Gate: adduct ratio within $\pm$2 SD
of cisplatin's value, OR ratio $\geq$ 2$\times$ for
candidates designed for \emph{enhanced} DNA binding.

\paragraph{Quality control.}
$^{195}$Pt blank $<$ 0.05\,pg; spike-recovery 85--115\%.
Each batch includes cisplatin-A2780 reference (gate:
18.5$\pm$5.0\,pg/\textmu g DNA; below triggers re-run).

\subsection{W.5 --- Timeline and cost}
\label{appendix:W-cost}

\paragraph{6-week plan.}
Week 1: building-block procurement (K$_2$[PtCl$_4$],
NH$_3$(aq), KI, AgNO$_3$, HCl, MTT, ICP-MS standards).
Week 2: cisplatin synthesis + characterisation.
Week 3: cell-line seeding + treatment (MCF-7, A549, A2780,
A2780cis).  Week 4: MTT 72\,h read-out + 4-plate Pt
uptake digestion.  Week 5: DNA extraction + ICP-MS on
DNA digests.  Week 6: data analysis + write-up of
measured-vs-predicted agreement.

\paragraph{Budget (¥ CNY).}
Reagents ¥ 5--8k, cell-line licence fees ¥ 4--6k,
consumables (RPMI/FBS/96-well plates) ¥ 3--5k,
ICP-MS instrument time ¥ 8--12k (3 sessions $\times$
$\sim$¥3k), labour (1 postdoc $\times$ 6 weeks) ¥ 30--40k.
\textbf{Total: ¥ 50--70k} (median ¥60k), well within the
appendix-W §W.5 budget envelope and comparable to a
PlatinAI-style 4-cell-line CRO quote ($\sim$¥80--100k at
2026 WuXi AppTec rates).

\paragraph{Honest framing.}
\textbf{None of the above has been executed.}  The
protocols are PROJECTED from the PlatinAI supplementary
material and standard Pt-cytoxicity bench literature
(Kauffman 1979 + Mosmann 1983 + Jamieson 1999).  The
data-side plumbing that the W.2--W.4 readouts would
feed (\texttt{wetlab\_protocol.py} +
\texttt{wetlab\_reward\_channel.py}) is SHIPPED at Tier-1
(\S\ref{sec:future:wetlab-validation}, 2026-09-17).
Tier-3 (collaborator outreach, IRB paperwork,
plate-reader integration) is the work this appendix plans.

\endinput